% 可从现有 paper/main.tex 用 input 引入；图像路径以该文件所在目录为基准。
\subsection{信号源数量与单源平均清除时长}
设信号源数量为$N$，全部清除总时长为$T$，单源平均时长为$\tau=T/N$。以下研究默认方法78在$10\le N\le16$内的条件均值，不将源数量视为决定单局时长的唯一因素。

总行走距离、换频数、测向数和失败清除数分别记为$L,S,M,F$，计时规则严格给出
\begin{equation}T=L/5+S+5M+3F+5N.\end{equation}
记实际发现点数为$m$，第$k$个源首次被发现的发现点序号为$J_k$。空频道在每一点检测，有源频道检测至首次发现，故
\begin{equation}M_{\mathrm{disc}}=m(20-N)+\sum_{k=1}^N J_k.\end{equation}
每一发现点的首次测量至多节省一次换频。记$d=M_{\mathrm{disc}}-S_{\mathrm{disc}}$，则$0\le d\le m$；其它检测和换频数记为$M_{\mathrm{loc}},S_{\mathrm{loc}}$。由此得到
\begin{equation}\begin{aligned}
\tau={}&\frac{L}{5N}+6m\left(\frac{20}{N}-1\right)+6\overline J+5\\
&+\frac{5M_{\mathrm{loc}}+S_{\mathrm{loc}}+3F-d}{N},\qquad \overline J=\frac1N\sum_kJ_k.
\end{aligned}\end{equation}
源数增加会减少空频道重复扫描，并分摊共同开销。$N<16$时，Q3与Q4分别需完成7个和21个发现点；$N=16$时，已知数量上界允许提前结束发现。

据此在有限数量范围内建立均值近似
\begin{equation}
\mathbb E[T\mid N]\approx A+BN+C\mathbf1_{\{N=16\}},\qquad
\mu(N):=\mathbb E[\tau\mid N]\approx B+\frac A N+\frac C N\mathbf1_{\{N=16\}}.
\end{equation}
原独立集每问200局用于事后关系建模；冻结公式后，新增种子2710001至2710200、每问200局作验证，不用新样本重估参数。总时长最小二乘拟合得到（单位：秒）
\begin{align}
\widehat\mu_3(N)&=60.01+\frac{2240.54}{N}-\frac{158.67}{N}\mathbf1_{\{N=16\}},\\
\widehat\mu_4(N)&=-31.60+\frac{6206.85}{N}-\frac{513.96}{N}\mathbf1_{\{N=16\}}.
\end{align}
HC3异方差稳健标准误给出两问$A$的95\%近似区间分别为$[2014.89,2466.19]$、$[5889.35,6524.36]$，均为正；$C$的相应区间均为负。因而对$N=10,\ldots,14$，模型预测的相邻变化为
\begin{equation}\widehat\mu(N+1)-\widehat\mu(N)=-\frac{\widehat A}{N(N+1)}<0.\end{equation}
从15到16还包含$\widehat C/16$的修正。$B$是减少空频道检测后的净增量系数，不是独立单源的操作费用，因此可以为负。

\begin{table}[H]\centering\small\caption{冻结的新200局验证，逐局计算误差}
\begin{tabular}{lrr}\toprule 指标&Q3&Q4\\\midrule
单源均时RMSE/秒&17.20&28.93\\
单源均时$R^2$&0.746&0.906\\
总时长$R^2$&0.203&0.364\\\bottomrule
\end{tabular}\end{table}
原400行结果逐行重现，800个剖析案例均通过计时及发现次数恒等式。公式主要解释每源平均时长的下降趋势；总时长解释力较低，空间位置、接收半径、朝向和误差仍造成较大变动。尤其16源分组存在可见预测误差，不能把公式当作单局精确预测或所有布局的单调性定理。结论限于含40米位置排斥规则的本地场景分布，不外推官方成绩或题设数量范围之外。

\begin{figure}[H]\centering
\includegraphics[width=\linewidth]{../补充分析/信号源数量与平均时长/figures/count_time_relation.pdf}
\caption{源数量与每源均时。阴影为拟合均值HC3近似95\%区间，误差棒为新样本分组均值区间；16处方点包含停止规则修正。}
\end{figure}
